% META: {"id": "1NSI-APT-RE-C1-CORRIGE", "chapitre": "1NSI-ALGO-PARCOURS-TRIS", "type_objet": "corrige", "exercice_ref": "1NSI-APT-RE-C1", "capacites": ["P-ALGO-01A"], "capacites_codes": ["C1"], "status": "needs_review"}
% BEGIN-VERIFY
% def bugue(tableau, valeur):
%     for i in range(len(tableau)):
%         if tableau[i] == valeur:
%             return i
%         else:
%             return None
% assert bugue([7, 3, 9], 3) is None
% assert bugue([7, 3, 9], 7) == 0
% assert bugue([7, 3, 9], 9) is None
% def correct(tableau, valeur):
%     for i in range(len(tableau)):
%         if tableau[i] == valeur:
%             return i
%     return None
% assert correct([7, 3, 9], 3) == 1
% assert correct([7, 3, 9], 9) == 2
% assert correct([7, 3, 9], 5) is None
% END-VERIFY

\begin{corrige}{1NSI-APT-RE-C1-EX1}

\textbf{1.} \lstinline{recherche_occurrence_bugue([7, 3, 9], 3)} renvoie \lstinline{None},
alors que $3$ se trouve bien en case $1$. Le résultat est \textbf{faux}.

\textbf{2.} Avec la valeur $7$, la fonction renvoie $0$, ce qui est correct. Un élève qui
ne testerait que ce cas conclurait que sa fonction marche. Elle ne donne le bon résultat que
lorsque la valeur cherchée occupe la \textbf{première} case — le seul cas où l'erreur est
invisible. Un test doit donc toujours viser aussi le milieu, la fin, et l'absence.

\textbf{3.} Un seul tour s'exécute, quel que soit le tableau. Au tour $i=0$, l'une des deux
branches est nécessairement empruntée : si \lstinline{tableau[0]} vaut la valeur cherchée,
le \lstinline{return i} sort ; sinon le \lstinline{return None} du \lstinline{else} sort. Or
un \lstinline{return} termine la fonction entière, pas seulement le tour. La boucle
n'atteint jamais l'indice $1$.

\textbf{4.} Le \lstinline{return None} doit être placé \textbf{après} la boucle, au niveau
d'indentation de la boucle elle-même :
\begin{python}
def recherche_occurrence(tableau, valeur):
    for i in range(len(tableau)):
        if tableau[i] == valeur:
            return i
    return None
\end{python}
La logique est asymétrique, et c'est la clé : pour affirmer « trouvé », une seule case
suffit — on peut donc sortir immédiatement. Pour affirmer « absent », il faut les avoir
\textbf{toutes} examinées — la conclusion ne peut donc être tirée qu'une fois la boucle
terminée.

\end{corrige}
